% Chapter 1 — Introduction (plain-language edition)

\chapter{Introduction}
\label{ch:intro}

\section{Why this project matters}
Schools and training centres increasingly teach online. That creates a simple question: \textit{what are learners actually doing}? Are they opening lessons? How long do they watch? Do they skip videos? How do they score on short quizzes?

If we store this information in spreadsheets or scattered files, reports soon disagree with reality. Names get duplicated, totals do not match, and nobody trusts the numbers.

A \textbf{database} fixes this by keeping \textbf{one official copy} of each fact. SQL makes it easy to ask questions like ``how many students finished Module 2?'' or ``who spent the least time on videos last week?'' The database also \textbf{blocks nonsense data}: for example, you cannot enroll someone in a course that does not exist, because the database enforces \textbf{foreign keys}.

\section{Problem we solved}
We needed a system that does three things at once:

\begin{enumerate}[label=(\alph*),leftmargin=*]
  \item \textbf{Day-to-day use.} Students register, join courses, watch content, and take quizzes. Instructors create courses and see basic reports. All of this must feel fast and reliable.
  \item \textbf{Detailed behaviour data.} We store \textbf{events} (play, pause, skip, etc.) and \textbf{watch time in seconds}, not just ``completed yes/no.'' That supports honest analytics later.
  \item \textbf{Reports without killing the server.} Counting and joining millions of raw rows on every button click is expensive. We therefore store some \textbf{pre-calculated summaries} (materialized views) that dashboards read quickly.
\end{enumerate}

\section{What we set out to achieve}
\begin{enumerate}[leftmargin=*]
  \item Design a \textbf{clear relational model}: separate tables for users, courses, enrollments, lessons, quizzes, attempts, and small ``lookup'' lists (status names, activity types, etc.).
  \item Put \textbf{business rules in the database} where possible: primary keys, foreign keys, checks, and uniqueness so the app cannot accidentally create duplicate enrollments or orphan rows.
  \item Build \textbf{analytics snapshots} (materialized views) keyed by student and course, so instructor dashboards stay responsive.
  \item Keep \textbf{history when we hide things}: soft deletes for users and courses, timestamps on important rows, and full quiz attempt + answer detail for auditing.
  \item Expose everything through a \textbf{REST API} (\texttt{/api/v1/...}) so the React app never talks SQL directly.
\end{enumerate}

\section{What this report covers (and what it does not)}
\textbf{In scope:} How tables relate to each other; why we chose UUIDs for people and integers for courses; how enrollments, activity, progress, and quizzes map to SQL; how materialized views and refresh functions support analytics; how API routes match database writes.

\textbf{Out of scope:} Machine learning grade prediction; native mobile apps; heavy research on learning psychology. Those could build \textit{on top of} this database later.

\section{How to read this report}
\begin{itemize}[leftmargin=*]
  \item Chapter~\ref{ch:overview} gives the big picture: browser, API, Supabase PostgreSQL.
  \item Chapter~\ref{ch:req} lists requirements and which tables satisfy them.
  \item Chapters~\ref{ch:conceptual} and~\ref{ch:logical} explain the entity model and normalization; placeholders show where your ERD and relational diagram go.
  \item Later chapters cover indexes, the full data dictionary, triggers, transactions, analytics queries, workflows, SQL examples, frontend data flow, testing, deployment, and conclusions.
\end{itemize}
